%% ma: L10-B13-0070
%% lop: 10
%% chuong: 5
%% bai: 13
%% dang: 10.13.10
%% loai: TL
%% muc: VD
%% dap_an: "-"
%% nguon: "chuyen-de-cac-so-dac-trung-cua-mau-so-lieu-khong-ghep-nhom-toan-10-chuong-trinh-moi.pdf (GV. Phan Nhật Linh) (P03), Bài 02 Dạng 2 Bài tập tự luận 7"
%% kien_thuc: "Lớp 10 Bài 13 (số trung bình, trung vị, tứ phân vị, mốt)"
%% trang_thai: chinh-thuc
%% ghi_chu: "Tự duyệt theo yêu cầu giáo viên 22/10/2026."
\begin{ex}
Hàm lượng natri (đơn vị: miligam) trong $100$ g của một số loại ngũ cốc được cho như sau:
\begin{center}
\begin{tabular}{|c|c|c|c|c|c|c|c|c|c|}\hline
$0$ & $340$ & $70$ & $140$ & $200$ & $180$ & $210$ & $150$ & $100$ & $130$\\\hline
$140$ & $180$ & $190$ & $160$ & $290$ & $50$ & $220$ & $180$ & $200$ & $210$\\\hline
\end{tabular}
\end{center}
Hãy tìm các tứ phân vị. Các tứ phân vị này cho ta thông tin gì?
\loigiai{Sắp xếp: $0$; $50$; $70$; $100$; $130$; $140$; $140$; $150$; $160$; $180$; $180$; $180$; $190$; $200$; $200$; $210$; $210$; $220$; $290$; $340$.

$n=20$: $Q_2=\dfrac{180+180}2=180$; $Q_1=\dfrac{130+140}2=135$; $Q_3=\dfrac{200+210}2=205$.

Khoảng từ $Q_1$ đến $Q_2$ dài $45$, còn từ $Q_2$ đến $Q_3$ chỉ dài $25$: số liệu bên phải $Q_2$ tập trung dày hơn, bên trái $Q_2$ phân bố thưa hơn.}
\end{ex}
